% !PRINTED 752
\weberpage{752}

%<<CONT
allm\"ahlich so v\"ollig beseitigt sind, dass sich der Beweis jetzt mit ganz elementaren Mitteln und auf die einfachste Weise f\"uhren l\"asst.

Wenn wir, wie schon fr\"uher, mit $\Pi(n)$ das Product
\[
\Pi(n) = 1.2.3\dots n
\]
bezeichnen, so wird, wenn $x$ eine beliebige Gr\"osse ist,
\weq{1}{\frac{x^n}{\Pi(n)}}
mit unendlich wachsendem $n$ sich der Grenze Null n\"ahern, und zwar st\"arker, als die Glieder einer fallenden geometrischen Reihe.

Denn ist $k$ eine ganze Zahl, die gr\"osser ist als der absolute Werth von $x$, so ist der absolute Werth von $x : h$ immer dann ein echter Bruch, wenn $h$ gleich oder gr\"osser als $k$ ist. Folglich ist
\[
\left|\frac{x^n}{\Pi(n)}\right| = \left|\frac{x^k}{\Pi(k)}\, \frac{x}{k+1}\, \frac{x}{k+2} \cdots \frac{x}{n}\right| < \left|\frac{x^k}{\Pi(k)}\right| \left|\frac{x}{k}\right|^{n-k}.
\]
% [sic] -- the strict inequality tacitly assumes $n>k$ and $x\ne0$; $x=0$ is the trivial equality and convergence case.
Hieraus folgt, dass die unendliche Reihe
\weq{2}{e^x = 1 + x + \frac{x^2}{\Pi(2)} + \frac{x^3}{\Pi(3)} + \cdots}
f\"ur alle Werthe von $x$ convergirt, und durch sie definiren wir die Exponentialfunction $e^x$, deren einfachste Eigenschaften wir hier aus der Analysis voraussetzen. Insbesondere gilt f\"ur zwei beliebige Werthe $x, y$ die Relation
\[
e^{x+y} = e^x e^y,
\]
aus der dann folgt, dass $e^n$ f\"ur ein ganzes positives $n$ die $n^{\text{te}}$ Potenz der Zahl
\weq{3}{e = 2 + \frac{1}{\Pi(2)} + \frac{1}{\Pi(3)} + \frac{1}{\Pi(4)} + \cdots = 2{,}718281828459\cdots}
ist.

Wenn nun $r$ irgend eine ganze positive Zahl bedeutet, so k\"onnen wir die Formel (2) so darstellen:
\weq{4}{\Pi(r)\, e^x = \Pi(r) + \frac{\Pi(r)}{\Pi(1)} x + \frac{\Pi(r)}{\Pi(2)} x^2 + \cdots + x^r + x^r\, U_r,}

\noindent\rule{4cm}{0.4pt}

Zahl\grqq{}. Sitzungsbericht der Berliner Akademie, 3. December 1885. Die Arbeiten von \gesperrt{Hilbert}, \gesperrt{Hurwitz} und \gesperrt{Gordan} finden sich alle drei in Band 43 der Mathematischen Annalen (1893), die beiden ersten auch in den G\"ottinger Nachrichten von 1893.
%CONT>>
